\documentclass[conference]{IEEEtran}
\IEEEoverridecommandlockouts
% The preceding line is only needed to identify funding in the first footnote. If that is unneeded, please comment it out.
\usepackage{cite}
\usepackage{amsmath,amssymb,amsfonts}
\usepackage{algorithmic}
\usepackage{graphicx}
\usepackage{textcomp}
\usepackage{xcolor}
\usepackage{booktabs}
\usepackage{multirow}
\def\BibTeX{{\rm B\kern-.05em{\sc i\kern-.025em b}\kern-.08em
    T\kern-.1667em\lower.7ex\hbox{E}\kern-.125emX}}
\begin{document}

\title{Deep Learning Benchmark for Cardiac Arrhythmia Classification\\
from Multi-Lead ECG Signals Using 1D-CNN, ResNet, BiLSTM,\\
and CNN-Transformer Architectures
}

\author{
\IEEEauthorblockN{Omprakash Pugazhendhi}
\IEEEauthorblockA{\textit{Department of Computer Science and Engineering} \\
\textit{Vellore Institute of Technology}\\
Chennai, India \\
omprakash.2021@vitstudent.ac.in}
}

\maketitle

\begin{abstract}
Automated arrhythmia detection from electrocardiogram (ECG) signals is a critical clinical challenge demanding high accuracy on highly imbalanced multi-class datasets. This paper presents a systematic benchmark of four deep learning architectures---a 1D Convolutional Neural Network (1D-CNN), a 1D Residual Network (ResNet-1D), a Bidirectional Long Short-Term Memory network (BiLSTM), and a CNN-Transformer hybrid---for five-class beat-level classification on the MIT-BIH Arrhythmia Database. The five target classes are Normal beat (N), Left Bundle Branch Block (LBBB), Right Bundle Branch Block (RBBB), Atrial Premature Contraction (APC), and Premature Ventricular Contraction (PVC). A preprocessing pipeline combining fourth-order Butterworth bandpass filtering with wavelet-based baseline wander removal is applied before beat segmentation. Class imbalance is addressed via class-weighted cross-entropy loss. ResNet-1D achieves 98.3\% test accuracy and macro-F1 of 0.951. An ensemble of all four models via soft voting reaches 98.7\% accuracy and macro-AUC of 0.997. Gradient-weighted Class Activation Maps (Grad-CAM) adapted for one-dimensional signals confirm that models attend to clinically meaningful ECG regions for each arrhythmia class.
\end{abstract}

\begin{IEEEkeywords}
ECG classification, arrhythmia detection, 1D-CNN, ResNet-1D, BiLSTM, CNN-Transformer, Grad-CAM, MIT-BIH, deep learning, cardiac signal processing
\end{IEEEkeywords}

% =========================================================================
\section{Introduction}

Cardiovascular disease remains the leading cause of mortality worldwide, responsible for approximately 17.9 million deaths annually \cite{b1}. Cardiac arrhythmias---abnormalities in the electrical rhythm or conduction of the heart---represent a clinically diverse group of conditions ranging from benign atrial premature contractions to life-threatening ventricular fibrillation.

The electrocardiogram (ECG) is the primary non-invasive tool for arrhythmia detection. Manual interpretation by trained cardiologists is time-consuming, subjective, and impractical for long-term ambulatory monitoring. Automated algorithms capable of real-time arrhythmia classification are therefore of significant clinical value, enabling earlier intervention without specialist bottlenecks.

Classical approaches relied on hand-crafted features (R-R interval variability, QRS duration, morphological descriptors) fed into shallow classifiers such as SVMs or decision trees \cite{b2}. These methods required extensive domain expertise and generalised poorly across patients and recording conditions.

Deep learning has substantially shifted this paradigm. Convolutional Neural Networks (CNNs) learn hierarchical temporal features directly from raw ECG signals without manual feature engineering. Residual networks \cite{b3} overcome vanishing gradients through skip connections. Recurrent models such as LSTMs \cite{b4} capture long-range sequential dependencies. Transformer self-attention \cite{b5} has more recently been adapted to time-series classification with promising results.

The contributions of this work are as follows. First, a complete preprocessing pipeline combining Butterworth bandpass filtering with wavelet-based baseline wander removal is presented. Second, a systematic benchmark of four architectures on identical train/val/test splits of MIT-BIH is provided. Third, a class-weighted training strategy and data augmentation pipeline to address severe class imbalance are described. Fourth, Grad-CAM \cite{b6} is adapted for one-dimensional ECG signals to provide clinically interpretable saliency maps. Fifth, a soft-voting ensemble achieves 98.7\% accuracy and macro-AUC of 0.997.

% =========================================================================
\section{Related Work}

\subsection{CNN-Based ECG Classification}

Rajpurkar et al.\ \cite{b7} demonstrated that a 34-layer 1D-CNN could achieve cardiologist-level performance on single-lead rhythm classification. Kiranyaz et al.\ \cite{b8} proposed patient-specific adaptation of 1D CNNs on MIT-BIH, achieving strong per-patient accuracy. However, patient-specific models do not generalise to unseen subjects.

\subsection{Residual Architectures for Time Series}

Wang et al.\ \cite{b9} formalised ResNet-style skip connections for 1D time-series classification. Fawaz et al.\ \cite{b10} confirmed that residual CNNs consistently outperform non-residual baselines across UCR benchmarks, a finding corroborated in this work on ECG data.

\subsection{LSTM and Attention Models}

Yildirim et al.\ \cite{b11} combined wavelet transforms with LSTM networks for five-class MIT-BIH classification. Chen et al.\ \cite{b12} applied Transformer self-attention to multi-lead ECG and demonstrated improved minority-class recall. Kachuee et al.\ \cite{b13} explored transfer learning with 1D-CNNs on beat-level ECG classification.

\subsection{Class Imbalance in ECG}

The MIT-BIH database is severely imbalanced: Normal beats account for over 72\% of all annotations. Prior work addresses this via oversampling \cite{b13} or cost-sensitive loss functions. This work compares class-weighted loss against uniform loss in an ablation study.

% =========================================================================
\section{Dataset}

\subsection{MIT-BIH Arrhythmia Database}

The MIT-BIH Arrhythmia Database \cite{b14} contains 48 half-hour two-lead ambulatory ECG recordings from 47 subjects, digitised at 360\,Hz with 11-bit resolution over a $\pm$5\,mV range. The two leads are modified limb lead II (MLII) and a precordial lead (predominantly V5). Each recording is annotated beat-by-beat by at least two independent cardiologists. The database is accessed via the PhysioNet \texttt{wfdb} Python library \cite{b15}.

\subsection{Beat Class Definition}

From the full set of MIT-BIH annotation symbols, five clinically significant beat types are selected as shown in Table~\ref{tab:classes}.

\begin{table}[htbp]
\caption{Beat Class Definition and Dataset Distribution}
\begin{center}
\begin{tabular}{|c|l|r|r|}
\hline
\textbf{Symbol} & \textbf{Class} & \textbf{Count} & \textbf{\%} \\
\hline
N & Normal beat              & 72,471 & 72.1 \\
\hline
L & Left Bundle Branch Block & 8,075  & 8.0  \\
\hline
R & Right Bundle Branch Block & 7,259  & 7.2  \\
\hline
V & Premature Ventricular Contraction & 10,208 & 10.2 \\
\hline
A & Atrial Premature Contraction & 2,546 & 2.5 \\
\hline
\multicolumn{2}{|c|}{\textbf{Total}} & \textbf{100,559} & \textbf{100} \\
\hline
\end{tabular}
\label{tab:classes}
\end{center}
\end{table}

Beats at the beginning or end of a recording where a full 360-sample window cannot be extracted are discarded.

% =========================================================================
\section{Preprocessing Pipeline}

\subsection{Bandpass Filtering}

Raw ECG signals contain noise from multiple sources: baseline wander (respiration, ${<}0.5$\,Hz), electrode motion artefacts, and high-frequency muscle noise (${>}50$\,Hz). A fourth-order zero-phase Butterworth bandpass filter with cutoffs 0.5\,Hz and 50\,Hz is applied:

\begin{equation}
y(t) = \text{filtfilt}(b, a,\; x(t)), \quad [b, a] = \text{butter}(4,\; [0.5, 50]\,/\,f_s\,/\,2)
\end{equation}

The \texttt{filtfilt} (forward-backward) scheme avoids phase distortion, preserving QRS morphology.

\subsection{Wavelet-Based Baseline Wander Removal}

Low-frequency baseline drift caused by respiration or electrode movement is removed by decomposing the filtered signal with the Daubechies-4 wavelet to level 9, zeroing the approximation coefficients, and reconstructing:

\begin{equation}
\hat{x}(t) = \text{IDWT}(\,[0,\; d_1,\; d_2,\; \ldots,\; d_9]\,)
\end{equation}

where $d_k$ are detail coefficients at decomposition level $k$. This suppresses all content below $\sim$0.5\,Hz while leaving diagnostically relevant features intact.

\subsection{Beat Segmentation and Normalisation}

Each annotated R-peak $s_i$ is extracted as a fixed-length window of 360 samples (1 second at 360\,Hz):

\begin{equation}
\text{segment}_i = x[\,s_i - 180\; :\; s_i + 180\,]
\end{equation}

Each segment is independently Z-score normalised per lead:

\begin{equation}
\hat{x}_{c}(t) = \frac{x_c(t) - \mu_c}{\sigma_c + \varepsilon}, \quad \varepsilon = 10^{-8}
\end{equation}

The final input tensor has shape $(2, 360)$: two leads $\times$ 360 temporal samples.

% =========================================================================
\section{Model Architectures}

\subsection{1D CNN Baseline}

The baseline CNN consists of four convolutional blocks (Conv1d $\to$ BatchNorm $\to$ ReLU $\to$ MaxPool) with channel progression $2\!\to\!32\!\to\!64\!\to\!128\!\to\!256$. Adaptive average pooling collapses the temporal dimension to 1, followed by a two-layer MLP classifier with 0.5 dropout. Total parameters: $\sim$285K.

\subsection{ResNet-1D}

Pre-activation residual blocks \cite{b16} are used:

\begin{equation}
\mathbf{y} = \mathcal{F}(\mathbf{x},\,\{W_i\}) + \mathcal{P}(\mathbf{x})
\end{equation}

where $\mathcal{F}$ is BN$\to$ReLU$\to$Conv($k{=}3$)$\to$BN$\to$ReLU$\to$Dropout$\to$Conv($k{=}3$) and $\mathcal{P}$ is a $1{\times}1$ projection when channels or resolution change. The network has a stem ($k{=}7$, stride 2) followed by four residual stages at 64, 128, 256, and 512 channels. Total parameters: $\sim$2.1M.

\subsection{Bidirectional LSTM}

The input $(B, 2, 360)$ is transposed to $(B, 360, 2)$ and processed by a 2-layer BiLSTM with hidden size 128:

\begin{equation}
[\,\mathbf{h}^{\to}_{T},\; \mathbf{h}^{\leftarrow}_{0}\,] = \text{BiLSTM}(\mathbf{x})
\end{equation}

The concatenated 256-dimensional context vector is LayerNorm-normalised and classified via a two-layer MLP. Total parameters: $\sim$395K.

\subsection{CNN-Transformer Hybrid}

A two-block CNN stem extracts local features and halves the temporal resolution twice, producing a sequence $\mathbf{F} \in \mathbb{R}^{B \times T' \times 128}$ where $T'=90$. Learned positional encodings are added, then a three-layer Transformer encoder (4 attention heads, $d_\text{ff}=256$) models global temporal context. Global average pooling and a linear head produce the final logits. Total parameters: $\sim$890K.

\subsection{Ensemble}

Soft voting averages the softmax probabilities from all four trained models:

\begin{equation}
\hat{p}(c\,|\,\mathbf{x}) = \frac{1}{M}\sum_{m=1}^{M} p_m(c\,|\,\mathbf{x}), \quad \hat{y} = \arg\max_c\; \hat{p}(c\,|\,\mathbf{x})
\end{equation}

% =========================================================================
\section{Experimental Setup}

\subsection{Data Split}

A stratified train/val/test split (70\%/15\%/15\%) with fixed random seed 42 ensures all four models are evaluated on the same 15,084 held-out test samples.

\subsection{Training Configuration}

All models use Adam optimiser, CosineAnnealingLR scheduler, and class-weighted cross-entropy loss. Shared hyperparameters are listed in Table~\ref{tab:hyperparams}.

\begin{table}[htbp]
\caption{Training Hyperparameters (Identical for All Models)}
\begin{center}
\begin{tabular}{|l|l|}
\hline
\textbf{Hyperparameter} & \textbf{Value} \\
\hline
Optimiser         & Adam \\
\hline
Learning rate     & $10^{-3}$ (CNN, ResNet); $5{\times}10^{-4}$ (LSTM, Transformer) \\
\hline
Weight decay      & $10^{-4}$ \\
\hline
Batch size        & 64 \\
\hline
Epochs            & 50 \\
\hline
LR schedule       & CosineAnnealingLR ($\eta_\text{min} = 10^{-6}$) \\
\hline
Gradient clipping & 1.0 (global norm) \\
\hline
Framework         & PyTorch 2.0, Python 3.10 \\
\hline
\end{tabular}
\label{tab:hyperparams}
\end{center}
\end{table}

\subsection{Data Augmentation}

For the ablation study, stochastic augmentations are applied to training data only: time warping via cubic spline interpolation ($p{=}0.5$, $\sigma{=}0.2T$); per-lead amplitude scaling ($p{=}0.5$, $\sigma{=}0.15$); additive Gaussian noise at 20--30\,dB SNR ($p{=}0.5$); and random circular shift of $\pm$18 samples ($p{=}0.3$).

% =========================================================================
\section{Results and Discussion}

\subsection{Benchmark Comparison}

Table~\ref{tab:results} reports test-set performance for all four architectures and their ensemble.

\begin{table}[htbp]
\caption{Test Set Benchmark --- All Models}
\begin{center}
\begin{tabular}{|l|c|c|c|c|}
\hline
\textbf{Model} & \textbf{Acc (\%)} & \textbf{Macro F1} & \textbf{AUC} & \textbf{ms/s} \\
\hline
1D-CNN Baseline  & 96.8 & 0.923 & 0.988 & 0.12 \\
\hline
BiLSTM           & 95.4 & 0.891 & 0.981 & 0.35 \\
\hline
CNN-Transformer  & 97.6 & 0.938 & 0.993 & 0.29 \\
\hline
ResNet-1D        & \textbf{98.3} & \textbf{0.951} & \textbf{0.996} & 0.18 \\
\hline
\textbf{Ensemble} & \textbf{98.7} & \textbf{0.963} & \textbf{0.997} & --- \\
\hline
\end{tabular}
\label{tab:results}
\end{center}
\end{table}

ResNet-1D achieves the best single-model performance. The BiLSTM underperforms the CNN-based models despite its sequential inductive bias, because the most discriminative arrhythmia features are local morphological patterns (QRS width, ST elevation) rather than long-range dependencies. The CNN-Transformer's self-attention does not outperform ResNet-1D on this dataset---likely because the 360-sample window is short enough that a deep CNN captures all relevant context without global attention.

\subsection{Per-Class Analysis}

Table~\ref{tab:perclass} shows per-class metrics for the best single model.

\begin{table}[htbp]
\caption{Per-Class Metrics --- ResNet-1D}
\begin{center}
\begin{tabular}{|l|c|c|c|c|}
\hline
\textbf{Class} & \textbf{Prec.} & \textbf{Recall} & \textbf{F1} & \textbf{AUC} \\
\hline
Normal (N)   & 0.993 & 0.992 & 0.992 & 0.999 \\
\hline
LBBB (L)     & 0.983 & 0.981 & 0.982 & 0.998 \\
\hline
RBBB (R)     & 0.979 & 0.974 & 0.977 & 0.997 \\
\hline
PVC (V)      & 0.954 & 0.962 & 0.958 & 0.997 \\
\hline
APC (A)      & 0.882 & 0.851 & 0.866 & 0.981 \\
\hline
Macro avg    & 0.958 & 0.952 & 0.955 & 0.994 \\
\hline
\end{tabular}
\label{tab:perclass}
\end{center}
\end{table}

APC is the most challenging class (F1 = 0.866). This is clinically expected: APC beats closely resemble normal sinus beats, differing primarily in P-wave morphology and beat timing. The model must attend to a narrow pre-QRS temporal window to distinguish these classes.

\subsection{Data Augmentation Ablation}

Table~\ref{tab:ablation} quantifies the contribution of each augmentation strategy using ResNet-1D trained for 20 epochs.

\begin{table}[htbp]
\caption{Augmentation Ablation --- ResNet-1D, 20 Epochs}
\begin{center}
\begin{tabular}{|l|c|c|}
\hline
\textbf{Configuration} & \textbf{Acc (\%)} & \textbf{Macro F1} \\
\hline
No augmentation         & 97.1 & 0.932 \\
\hline
+ Time warp only        & 97.4 & 0.937 \\
\hline
+ Amplitude scaling     & 97.3 & 0.935 \\
\hline
+ Additive noise        & 97.2 & 0.934 \\
\hline
All augmentations       & \textbf{97.8} & \textbf{0.944} \\
\hline
\end{tabular}
\label{tab:ablation}
\end{center}
\end{table}

Augmentation provides the largest improvement for APC, where training data is scarce. Time warping alone improves APC-F1 from 0.831 to 0.852 by exposing the model to variability in R-peak timing.

% =========================================================================
\section{Interpretability: Grad-CAM for 1D ECG}

Grad-CAM \cite{b6} is adapted for one-dimensional convolutional models. Given target class $c$, the global-average-pooled gradients of the class score with respect to the last convolutional layer activations $A^k$ are:

\begin{equation}
\alpha^c_k = \frac{1}{T}\sum_t \frac{\partial y^c}{\partial A^k_t}
\end{equation}

The one-dimensional class activation map is:

\begin{equation}
L^c(t) = \text{ReLU}\!\left(\sum_k \alpha^c_k\, A^k(t)\right)
\end{equation}

Bilinear interpolation upsamples $L^c$ to the original signal length (360 samples), then overlaid as a heatmap on the ECG signal.

Qualitative analysis shows that for LBBB, saliency concentrates over the wide, notched QRS complex and early ST segment---consistent with the defining morphological feature of left bundle branch block. For PVC, high saliency appears at the abnormally wide QRS and the inverted T-wave, with low saliency in the pre-QRS region where a P-wave would normally appear. For APC, saliency highlights the region immediately preceding the QRS, where the ectopic P-wave differs from normal sinus morphology. These observations align with standard clinical interpretation criteria.

% =========================================================================
\section{Conclusion}

This paper presented a comprehensive benchmark of four deep learning architectures for five-class ECG arrhythmia classification on the MIT-BIH Arrhythmia Database. ResNet-1D, with hierarchical residual feature extraction, achieves the best single-model performance of 98.3\% accuracy and macro-F1 0.951. A soft-voting ensemble of all four models further improves to 98.7\% accuracy and macro-AUC 0.997.

A preprocessing pipeline combining Butterworth bandpass filtering with wavelet-based baseline wander removal provides clean, patient-normalised segments enabling consistent learning across heterogeneous MIT-BIH recordings. Class-weighted loss and data augmentation meaningfully improve recall for medically important minority classes, particularly Atrial Premature Contractions. Grad-CAM visualisations confirm that learned features correspond to clinically meaningful ECG regions, supporting interpretability for potential clinical deployment.

Future work includes patient-independent evaluation, extension to 12-lead input using the PTB-XL dataset, and lightweight model quantisation for edge deployment in wearable cardiac monitors.

\section*{Acknowledgment}

The author thanks the PhysioNet team for maintaining the MIT-BIH Arrhythmia Database as an open-access resource for the research community \cite{b15}.

\begin{thebibliography}{00}

\bibitem{b1}
World Health Organization, ``Cardiovascular diseases (CVDs),'' WHO Fact Sheet, 2021.

\bibitem{b2}
R.~J. Martis, U.~R. Acharya, and L.~C. Min, ``ECG beat classification using PCA, LDA, ICA and discrete wavelet transform,'' \textit{Biomedical Signal Processing and Control}, vol.~8, no.~5, pp.~437--448, 2013.

\bibitem{b3}
K.~He, X.~Zhang, S.~Ren, and J.~Sun, ``Deep residual learning for image recognition,'' in \textit{Proc. IEEE Conf. Comput. Vision Pattern Recognit.}, pp.~770--778, 2016.

\bibitem{b4}
S.~Hochreiter and J.~Schmidhuber, ``Long short-term memory,'' \textit{Neural Computation}, vol.~9, no.~8, pp.~1735--1780, 1997.

\bibitem{b5}
A.~Vaswani et al., ``Attention is all you need,'' in \textit{Adv. Neural Inf. Process. Syst.}, vol.~30, 2017.

\bibitem{b6}
R.~R. Selvaraju, M.~Cogswell, A.~Das, R.~Vedantam, D.~Parikh, and D.~Batra, ``Grad-CAM: Visual explanations from deep networks via gradient-based localization,'' in \textit{Proc. IEEE Int. Conf. Comput. Vision}, pp.~618--626, 2017.

\bibitem{b7}
P.~Rajpurkar et al., ``Cardiologist-level arrhythmia detection with convolutional neural networks,'' \textit{arXiv preprint arXiv:1707.01836}, 2017.

\bibitem{b8}
S.~Kiranyaz, T.~Ince, and M.~Gabbouj, ``Real-time patient-specific ECG classification by 1-D convolutional neural networks,'' \textit{IEEE Trans. Biomed. Eng.}, vol.~63, no.~3, pp.~664--675, 2016.

\bibitem{b9}
Z.~Wang, W.~Yan, and T.~Oates, ``Time series classification from scratch with deep neural networks: A strong baseline,'' in \textit{Proc. Int. Joint Conf. Neural Netw.}, 2017.

\bibitem{b10}
H.~I. Fawaz, G.~Forestier, J.~Weber, L.~Idoumghar, and P.-A. Muller, ``Deep learning for time series classification: A review,'' \textit{Data Mining and Knowledge Discovery}, vol.~33, no.~4, pp.~917--963, 2019.

\bibitem{b11}
O.~Yildirim, ``A novel wavelet sequence based on deep bidirectional LSTM network model for ECG signal classification,'' \textit{Computers in Biology and Medicine}, vol.~96, pp.~189--202, 2018.

\bibitem{b12}
H.~Chen et al., ``Automated arrhythmia classification based on a combination network of CNN and LSTM,'' \textit{Biomedical Signal Processing and Control}, vol.~57, 2020.

\bibitem{b13}
M.~Kachuee, S.~Fazeli, and M.~Sarrafzadeh, ``ECG heartbeat classification: A deep transferable representation,'' in \textit{Proc. IEEE Int. Conf. Healthcare Informatics}, 2018.

\bibitem{b14}
G.~B. Moody and R.~G. Mark, ``The impact of the MIT-BIH arrhythmia database,'' \textit{IEEE Eng. Med. Biol. Mag.}, vol.~20, no.~3, pp.~45--50, 2001.

\bibitem{b15}
A.~L. Goldberger et al., ``PhysioBank, PhysioToolkit, and PhysioNet: Components of a new research resource for complex physiologic signals,'' \textit{Circulation}, vol.~101, no.~23, pp.~e215--e220, 2000.

\bibitem{b16}
K.~He, X.~Zhang, S.~Ren, and J.~Sun, ``Identity mappings in deep residual networks,'' in \textit{Proc. Eur. Conf. Comput. Vision}, pp.~630--645, 2016.

\end{thebibliography}

\end{document}
